\clearpage
\SourceLeaf{237}%
\begin{center}
\Large\textsc{[Lübsen's theory of parallels.]}
\end{center}


\begin{center}
\emph{[\textup{\textsc{Schumacher}} to \textup{\textsc{Gauss}}. Altona, 31 December 1835.]}
\end{center}

\{\mbox{.\ .\ .\ .\ .\ .\ .\ .\ .} \textsc{Lübsen} is tormenting me with his theory of parallels
(Nemesis is probably stepping in here to repay me for my own
attempts). I enclose his proof. I pointed out to him
in vain that shifting lines is better done with a
parallel ruler than in theory. He still holds to his
conviction that he has found the geometrical proof.
Since he is otherwise a sensible and
modest man, it might be a good thing if you could disabuse him
in a few words. The trouble must be that I cannot make
his error as clear to him as I see it
myself. \mbox{.\ .\ .\ .\ .\ .}\}


\begin{center}
\textsc{Gauss} to \textsc{Schumacher}. Göttingen, 2 January 1836.
\end{center}

\mbox{.\ .\ .\ .\ .\ .\ .\ .} The crux of Mr \textsc{Lübsen}'s error is that
Euclid's geometry can be false without his construction having to yield,
in
